% 03_framework.tex — Rail Paper: The SolveEverything Framework
% ============================================================

\section{The SolveEverything Framework for Imaging}
\label{sec:framework}

The SolveEverything.org framework~\cite{wissner2026solveeverything} identifies a
recurring pattern across domains that have undergone AI-driven
transformation: the decisive enabler is never the model alone but the
\emph{infrastructure} that makes models measurable, comparable, and
improvable.
This section maps the framework's two core structures---the four stages
of revolution and the nine-layer Industrial Intelligence Stack---to
computational imaging, producing a concrete roadmap for moving the field
from its current pre-industrial state to systematic, scalable progress.

% ------------------------------------------------------------------
%  4 stages of revolution
% ------------------------------------------------------------------

\subsection{Four Stages of Revolution}
\label{sec:stages}

SolveEverything identifies four stages through which a domain passes
on its way from artisanal practice to industrial abundance.
We map each stage to computational imaging:

\paragraph{Stage 1: Legibility.}
\textit{Can you measure the problem?}
In imaging terms: can you define what ``reconstruction quality'' means
in a way that all practitioners agree on?
Currently, the answer is no.
As documented in \cref{sec:fragmentation}, the field uses
incommensurable metrics, datasets, noise models, and mismatch
definitions.
The first prerequisite for progress is a shared measurement framework.

\paragraph{Stage 2: Harnessing.}
\textit{Can you capture the energy?}
In imaging terms: can you standardize evaluation so that improvements
measured by one group are meaningful to all others?
This requires not only agreed-upon metrics but a common protocol for
applying them---including standardized forward operators, mismatch
injection procedures, and reporting formats.
The 4-scenario protocol and OperatorGraph IR described in
\cref{sec:architecture} are designed to achieve exactly this.

\paragraph{Stage 3: Institutionalization.}
\textit{Can you scale it?}
In imaging terms: can you automate calibration and evaluation so that
they do not require per-system, per-lab, per-paper manual effort?
This is the stage at which LIP-Arena and the Red Team module operate:
automated, prospective evaluation that runs continuously and at scale,
with institutional governance to ensure integrity.

\paragraph{Stage 4: Abundance.}
\textit{Is it a commodity?}
In imaging terms: calibration as a utility.
Every imaging system ships with a self-calibrating operator model that
continuously updates itself, and reconstruction quality is guaranteed by
infrastructure rather than heroic engineering.
This is the end state that PWM is designed to enable, though the field
is far from reaching it today.

% ------------------------------------------------------------------
%  Rails vs Trains
% ------------------------------------------------------------------

\subsection{Rails vs.\ Trains}
\label{sec:rails_trains}

The SolveEverything framework draws a fundamental distinction between
\emph{rails} and \emph{trains}:

\begin{quote}
\textit{``The durable, compounding value\ldots will not be found in
owning any single AI model.
Models are the `trains' that will eventually all look the same.
The real value lies in owning the `rails'.''}
\hfill---SolveEverything.org~\cite{wissner2026solveeverything}
\end{quote}

In computational imaging, the trains are reconstruction
algorithms---GAP-TV, MST-L, HDNet, and their successors.
The history of deep learning in other domains (\eg, image
classification, natural language processing) strongly suggests that
reconstruction algorithms will converge: architectures will
commoditize, pre-trained foundation models will emerge, and the
marginal improvement from any single new method will shrink
asymptotically.

The rails are the infrastructure components that remain valuable
regardless of which solver is running:
\begin{itemize}
  \item The OperatorGraph IR that represents forward models in a
        universal, machine-readable format.
  \item The 4-scenario evaluation protocol that separates solver quality
        from operator fidelity.
  \item The LIP-Arena that provides prospective, anti-Goodhart
        evaluation.
  \item The calibration pipelines that correct operators before
        reconstruction.
  \item The governance structures that ensure reproducibility and
        accountability.
\end{itemize}

PWM is designed entirely as a rail.
It does not include a reconstruction algorithm; it includes the
infrastructure that makes reconstruction algorithms measurable.

% ------------------------------------------------------------------
%  9-layer Industrial Intelligence Stack
% ------------------------------------------------------------------

\subsection{The Industrial Intelligence Stack}
\label{sec:stack}

SolveEverything defines a nine-layer stack that organizes the
components needed for industrializing AI in any domain.
\Cref{tab:stack_mapping}
provides the detailed mapping of each layer to the corresponding PWM component.

\begin{table}[t]
\centering
\caption{%
  \textbf{SolveEverything Industrial Intelligence Stack mapped to PWM.}
  Each layer of the generic stack is instantiated with the corresponding
  PWM component and its concrete implementation.%
}
\label{tab:stack_mapping}
\smallskip
\renewcommand{\arraystretch}{1.15}
\begin{tabular}{@{}clll@{}}
\toprule
\textbf{Layer} & \textbf{Generic Name}   & \textbf{PWM Component}         & \textbf{Implementation} \\
\midrule
1 & Purpose \& Payoff       & Recovery targets               & Recovery ratio $\geq 0.80$; oracle gap $\leq 2$\,dB \\
2 & Task Taxonomy           & OperatorGraph IR                & 64 modalities, 89 templates \\
3 & Observability           & Diagnostics                     & RunBundle, DR-IS, TriadReport \\
4 & Targeting System        & LIP-Arena                       & Commit-Measure-Score protocol \\
5 & Model Layer             & Reconstruction methods          & GAP-TV, MST-L, HDNet, \etc{} \\
6 & Actuation               & Operator correction             & Calibrated $\hat{\mathbf{H}}$ fed to solvers \\
7 & Verification            & Red Team module                 & 2{,}900+ adversarial tests \\
8 & Governance              & Outcome-based ranking           & Compute escrow, open reports \\
9 & Distribution            & Universal protocol              & OperatorGraph IR as interchange format \\
\bottomrule
\end{tabular}
\end{table}

We briefly describe the non-obvious layers:

\paragraph{Layer 1: Purpose \& Payoff.}
Before building infrastructure, one must define what success looks like.
For PWM, we set two concrete thresholds: a \emph{recovery ratio}
(corrected PSNR divided by ideal PSNR) of at least 0.80, meaning that
calibration recovers at least 80\% of the oracle performance; and an
\emph{oracle gap} (ideal PSNR minus corrected PSNR) of at most 2\,dB.
These thresholds operationalize the concept of ``good enough
calibration'' and provide a pass/fail criterion for the entire system.

\paragraph{Layer 3: Observability.}
A system that cannot be observed cannot be improved.
PWM's observability layer centers on three instruments:
(i)~\textbf{RunBundle}, the immutable record of every evaluation run,
including all inputs, outputs, hyperparameters, and random seeds;
(ii)~\textbf{DR-IS} (Degradation-Recovery Importance Sampling), a
diagnostic that identifies which mismatch parameters contribute most to
reconstruction loss; and
(iii)~\textbf{TriadReport}, a standardized 3-panel report showing
ideal, mismatched, and corrected reconstructions side by side for every
test scene.

\paragraph{Layer 4: Targeting System.}
The LIP-Arena (described in detail in \cref{sec:architecture})
functions as the targeting system: it identifies which problems
matter most and directs community effort toward them.
By publishing challenge tracks organized around specific mismatch
types and modalities, LIP-Arena ensures that research effort is
allocated to the highest-impact problems rather than the most
convenient benchmarks.

\paragraph{Layer 7: Verification.}
The Red Team module serves as the adversarial verification layer.
With over 2{,}900 pre-designed test scenarios spanning novel mismatch
types, compound perturbations, out-of-family physics, distribution
shift, and compute traps, it probes submitted methods for failure modes
that retrospective evaluation would never reveal.

\paragraph{Layer 8: Governance.}
Outcome-based ranking means that methods are ranked by their
\emph{actual deployed performance} (including mismatch and calibration),
not by their ideal-condition leaderboard numbers.
Compute escrow ensures that all methods are evaluated under comparable
computational budgets, preventing the conflation of algorithmic
improvement with hardware scaling.

% ------------------------------------------------------------------
%  L0–L5 maturity ladder
% ------------------------------------------------------------------

\subsection{The L0--L5 Maturity Ladder}
\label{sec:maturity}

SolveEverything defines six maturity levels for any domain's
infrastructure.
\Cref{tab:maturity} maps these levels to computational imaging.

\begin{table}[t]
\centering
\caption{%
  \textbf{L0--L5 maturity ladder for computational imaging.}
  Current state: transitioning from L1 to L2.
  PWM provides the infrastructure needed to reach L3 and beyond.%
}
\label{tab:maturity}
\smallskip
\renewcommand{\arraystretch}{1.15}
\begin{tabular}{@{}clp{7.8cm}@{}}
\toprule
\textbf{Level} & \textbf{Name}       & \textbf{Imaging Definition} \\
\midrule
L0 & The Muddle      & No agreement on metrics, datasets, or mismatch
                        definitions.  Every lab runs its own evaluation.
                        Results are mutually incomparable.
                        \emph{(Pre-PWM status quo.)} \\
L1 & Measurable      & Clear metrics exist.  The 4-scenario protocol
                        defines what to measure and how.  Mismatch is
                        acknowledged as a first-class evaluation axis.
                        \emph{(InverseNet contribution.)} \\
L2 & Repeatable      & Standard operating procedures (SOPs) are
                        documented and published.  Any lab can reproduce
                        any other lab's evaluation using the same
                        OperatorGraph templates, datasets, and scoring
                        code. \\
L3 & Automated       & 80\% of calibration and evaluation runs without
                        human intervention.  LIP-Arena runs autonomously.
                        New modalities can be on-boarded via templated
                        OperatorGraph definitions. \\
L4 & Industrialized  & Calibration as a service.  Hardware vendors ship
                        imaging systems with PWM-compatible operator
                        models.  Evaluation is continuous and
                        infrastructure-grade. \\
L5 & Commoditized    & Universal self-calibration.  Every imaging system
                        continuously estimates and corrects its own
                        forward operator.  Reconstruction quality is
                        guaranteed by infrastructure. \\
\bottomrule
\end{tabular}
\end{table}

\paragraph{Current position: L1 $\to$ L2.}
The InverseNet benchmark~\cite{yang2026inversenet} established L1 by
demonstrating that mismatch is measurable and defining a 3-scenario
protocol (Ideal, Mismatch, Oracle). This work extends InverseNet's framework to a 4-scenario protocol: InverseNet's Ideal and Mismatch map to our Scenarios~I and~II; InverseNet's Oracle (true operator on mismatched data) maps to our Scenario~IV (Oracle Mask); and Scenario~III (Corrected) is a new contribution. The protocol was validated across three modalities (CASSI, CACTI, SPC).
PWM extends this to a 4-scenario protocol---adding a Corrected scenario
that measures calibration effectiveness---across 64 modalities and
provides the OperatorGraph IR, RunBundle infrastructure, and LIP-Arena
governance needed to reach L2 (repeatability) and begin the transition
to L3 (automation).
The gap from L2 to L3 is primarily an engineering challenge:
automating operator calibration, scaling the evaluation pipeline, and
building the institutional capacity to run LIP-Arena rounds
continuously.

\paragraph{The critical transition: L2 $\to$ L3.}
The most impactful transition in the maturity ladder is from L2
(repeatable, but manual) to L3 (automated).
History suggests that this is where transformative progress occurs: it
is the transition from artisanal craft to industrial process, the point
at which the rate of improvement shifts from human-limited to
infrastructure-limited.
In protein folding, the analogous transition was CASP's evolution from
a biennial manual assessment to an automated, continuous evaluation
pipeline---a transition that directly enabled the AlphaFold
breakthrough~\cite{jumper2021alphafold}.
PWM's architecture is designed to make this transition possible for
computational imaging.

\subsection{From Framework to Architecture}
\label{sec:framework_to_arch}

The SolveEverything framework provides the \emph{what}: a clear
enumeration of the infrastructure components needed and the maturity
milestones that mark progress.
The next section provides the \emph{how}: the concrete architecture of
PWM, including the OperatorGraph IR specification, the 4-scenario
evaluation protocol, the LIP-Arena competition design, and the Red Team
adversarial evaluation module.
